\documentclass[12pt]{article}
\usepackage[utf8]{inputenc}
\usepackage[T1]{fontenc}
\usepackage[margin=1in]{geometry}
\usepackage{setspace}
\usepackage{xcolor}
\usepackage[colorlinks=true, urlcolor=blue!60!black, linkcolor=blue!60!black]{hyperref}
\setlength{\parindent}{2em}
\setlength{\parskip}{1.25em}
\renewcommand{\baselinestretch}{1.0}
\begin{document}
\begin{center}
{\LARGE \textit{Dobbs v. Jackson Women's Health Organization}, 597 U.S. 215 (2022)\par}
\vspace{0.8em}
{\large Case Brief\par}
\end{center}
\vspace{1em}

\section{Memory Jogger}

In \href{https://www.courtlistener.com/opinion/6481357/dobbs-v-jackson-womens-health-organization/}{\textit{Dobbs v. Jackson Women's Health Organization}}, the Supreme Court overruled \textit{Roe v. Wade} and \textit{Planned Parenthood of Southeastern Pennsylvania v. Casey}, holding that the Constitution confers no right to abortion and returning the authority to regulate or prohibit abortion to the people and their elected representatives.

\section{Detailed Case Facts}

In 2018, the Mississippi Legislature enacted the Gestational Age Act, codified at Miss.\ Code Ann.\ \S~41-41-191, which prohibited any person from intentionally or knowingly performing or inducing an abortion if the probable gestational age of the unborn human being was determined to be greater than 15 weeks (\href{https://www.law.cornell.edu/supremecourt/text/19-1392}{Cornell LII}; \href{https://www.oyez.org/cases/2021/19-1392}{Oyez}). The statute provided narrow exceptions only for a medical emergency or a severe fetal abnormality, and included no exceptions for pregnancies resulting from rape or incest (\href{https://constitutioncenter.org/the-constitution/supreme-court-case-library/dobbs-v-jackson-womens-health-organization}{National Constitution Center}). Legislative findings accompanying the Act asserted that by 15 weeks a fetus has developed certain physical features and that the state possessed legitimate interests in protecting maternal health, preserving fetal life, and preventing procedures the legislature characterized as barbaric (\href{https://www.supremecourt.gov/opinions/21pdf/19-1392_6j37.pdf}{Slip Opinion}).

At the time of enactment, Jackson Women's Health Organization (``JWHO'') was the only licensed abortion clinic operating in Mississippi, performing surgical abortions up to 16 weeks' gestation (\href{https://www.oyez.org/cases/2021/19-1392}{Oyez}). Within a day of the law's passage, JWHO and one of its physicians, Dr.\ Sacheen Carr-Ellis, filed suit in federal district court against Mississippi State Health Officer Thomas E. Dobbs and Kenneth Cleveland, Executive Director of the Mississippi State Board of Medical Licensure, seeking declaratory and injunctive relief (\href{https://www.courtlistener.com/opinion/6481357/dobbs-v-jackson-womens-health-organization/}{CourtListener}).

The dispute squarely implicated the Fourteenth Amendment's Due Process Clause---which prohibits any State from depriving ``any person of life, liberty, or property, without due process of law''---and nearly half a century of Supreme Court substantive due process jurisprudence. In \textit{Roe v. Wade}, 410 U.S. 113 (1973), the Court had recognized a constitutional right to abortion under the Fourteenth Amendment's concept of personal liberty, and in \textit{Planned Parenthood of Southeastern Pennsylvania v. Casey}, 505 U.S. 833 (1992), a controlling plurality had reaffirmed \textit{Roe}'s ``central holding'' that a State may not prohibit abortion prior to fetal viability (typically around 23 to 24 weeks' gestation) or impose an ``undue burden'' on pre-viability access (\href{https://constitutioncenter.org/the-constitution/supreme-court-case-library/dobbs-v-jackson-womens-health-organization}{National Constitution Center}).

Politically and institutionally, \textit{Dobbs} arrived after a decisive ideological shift on the Supreme Court following the appointments of Justices Neil Gorsuch, Brett Kavanaugh, and Amy Coney Barrett, giving conservative originalists and textualists a 6--3 majority. Roughly two dozen states had enacted or were advancing pre-viability restrictions, and thirteen states had ``trigger laws'' poised to ban nearly all abortions immediately upon \textit{Roe}'s overruling. In an unprecedented breach of Court confidentiality, a draft of Justice Alito's majority opinion leaked to the press in May 2022, weeks before the final decision issued on June 24, 2022 (\href{https://www.supremecourt.gov/opinions/21pdf/597us1r58_gebh.pdf}{597 U.S. 215 Page Proof}).

\section{Procedural History}

\textbf{U.S. District Court for the Southern District of Mississippi:} JWHO filed suit seeking an emergency temporary restraining order (``TRO'') and permanent injunction against enforcement of the Gestational Age Act (\href{https://www.oyez.org/cases/2021/19-1392}{Oyez}). Judge Carlton W. Reeves granted a TRO and, following discovery, granted summary judgment in favor of JWHO and permanently enjoined the Act in \textit{Jackson Women's Health Org. v. Currier}, 349 F. Supp. 3d 536 (S.D. Miss. 2018) (\href{https://www.law.cornell.edu/supremecourt/text/19-1392}{Cornell LII}). Judge Reeves reasoned that ``viability marks the earliest point at which the State's interest in fetal life is constitutionally adequate to justify a legislative ban on nontherapeutic abortions'' and that 15 weeks' gestational age is indisputably ``prior to viability'' (\href{https://www.law.cornell.edu/supremecourt/text/19-1392}{Cornell LII}).

\textbf{U.S. Court of Appeals for the Fifth Circuit:} Mississippi appealed. In December 2019, a unanimous three-judge panel affirmed the permanent injunction in \textit{Jackson Women's Health Org. v. Dobbs}, 945 F.3d 265 (5th Cir. 2019), with Senior Circuit Judge Patrick Higginbotham writing that binding Supreme Court precedent from \textit{Roe} and \textit{Casey} categorically prohibited pre-viability bans (\href{https://www.law.cornell.edu/supremecourt/text/19-1392}{Cornell LII}).

\textbf{U.S. Supreme Court:} Mississippi petitioned for a writ of certiorari. On May 17, 2021, the Supreme Court granted certiorari limited to Question 1: whether ``all pre-viability prohibitions on elective abortions are unconstitutional'' (\href{https://www.oyez.org/cases/2021/19-1392}{Oyez}; \href{https://www.law.cornell.edu/supremecourt/text/19-1392}{Cornell LII}). Oral argument took place on December 1, 2021, featuring Mississippi Solicitor General Scott G. Stewart for petitioners, Julie Rikelman for respondents, and U.S. Solicitor General Elizabeth B. Prelogar as amicus curiae supporting respondents (\href{https://www.oyez.org/cases/2021/19-1392}{Oyez}). On June 24, 2022, the Court reversed the Fifth Circuit's judgment and remanded for further proceedings (\href{https://www.courtlistener.com/opinion/6481357/dobbs-v-jackson-womens-health-organization/}{CourtListener}). On remand, the Fifth Circuit returned the case to the District Court, which vacated its injunctions in September 2022 and entered final judgment for the state defendants after JWHO closed its Mississippi clinic (\href{https://clearinghouse.net/case/43153/}{Civil Rights Litigation Clearinghouse}).

\section{Judicial Votes}

The Court voted \textbf{6--3} to reverse the Fifth Circuit and uphold Mississippi's 15-week ban, and \textbf{5--4} (5--1--3) to overrule \textit{Roe v. Wade} and \textit{Planned Parenthood of Southeastern Pennsylvania v. Casey} in full (\href{https://constitutioncenter.org/the-constitution/supreme-court-case-library/dobbs-v-jackson-womens-health-organization}{National Constitution Center}).

\begin{itemize}
  \item \textbf{Majority (5):} Justice Samuel A. Alito, Jr.\ (author), joined by Justices Clarence Thomas, Neil M. Gorsuch, Brett M. Kavanaugh, and Amy Coney Barrett.
  \item \textbf{Concurrences Joining the Majority (2):} Justice Clarence Thomas (filed a separate concurring opinion); Justice Brett M. Kavanaugh (filed a separate concurring opinion).
  \item \textbf{Concurrence in the Judgment Only (1):} Chief Justice John G. Roberts, Jr.\ (agreed to uphold Mississippi's 15-week ban and discard the viability line, but declined to join the majority in overruling \textit{Roe} and \textit{Casey} entirely).
  \item \textbf{Joint Dissent (3):} Justices Stephen G. Breyer, Sonia Sotomayor, and Elena Kagan (co-authored a single joint dissenting opinion).
\end{itemize}

\section{Holding}

\textbf{Judgment:} Reversed and remanded. Mississippi's Gestational Age Act is constitutional under rational-basis review (\href{https://www.supremecourt.gov/opinions/21pdf/19-1392_6j37.pdf}{Slip Opinion}).

\textbf{Rule Announced:} The U.S. Constitution does not confer a right to abortion; \textit{Roe v. Wade} and \textit{Planned Parenthood of Southeastern Pennsylvania v. Casey} are overruled. Under the Fourteenth Amendment's Due Process Clause, an unenumerated right is protected only if it is ``deeply rooted in this Nation's history and tradition'' and ``implicit in the concept of ordered liberty'' under \textit{Washington v. Glucksberg}, 521 U.S. 702, 721 (1997). Because abortion does not meet that historical standard, and because abortion regulation is not a sex-based classification triggering heightened scrutiny under the Equal Protection Clause, state abortion laws are entitled to a ``strong presumption of validity'' and must be sustained under rational-basis review if there is a rational basis on which the legislature could have thought the law would serve legitimate state interests (\href{https://constitutioncenter.org/the-constitution/supreme-court-case-library/dobbs-v-jackson-womens-health-organization}{National Constitution Center}; \href{https://www.supremecourt.gov/opinions/21pdf/597us1r58_gebh.pdf}{597 U.S. 215 Page Proof}).

\textbf{Holding vs.\ Dicta:} The core holdings are that (1) neither the Due Process Clause nor the Equal Protection Clause of the Fourteenth Amendment protects a right to abortion, (2) \textit{Roe} and \textit{Casey} are overruled, and (3) rational-basis review governs constitutional challenges to abortion regulations. By contrast, the majority's repeated assurances that ``[n]othing in this opinion should be understood to cast doubt on precedents that do not concern abortion'' (such as \textit{Griswold v. Connecticut}, \textit{Lawrence v. Texas}, and \textit{Obergefell v. Hodges}) are formally \textit{dicta} because those precedents were not before the Court (\href{https://constitutioncenter.org/the-constitution/supreme-court-case-library/dobbs-v-jackson-womens-health-organization}{National Constitution Center}). Likewise, Justice Kavanaugh's statements regarding the constitutional right to interstate travel and the Ex Post Facto Clause, and Justice Thomas's call to reconsider all substantive due process decisions, are non-binding separate commentary.

\section{Analysis of Opinions}

\subsection{Majority Opinion (Justice Alito)}

Justice Alito structured the Court's opinion around three pillars: constitutional text and history, the distinction between abortion and other unenumerated rights, and a five-factor \textit{stare decisis} analysis (\href{https://www.law.cornell.edu/supremecourt/text/19-1392}{Cornell LII}). First, after briefly rejecting an Equal Protection Clause rationale as foreclosed by \textit{Geduldig v. Aiello} (holding that regulating a procedure unique to one sex is not an invidious sex-based classification absent pretext), Alito turned to the Fourteenth Amendment's Due Process Clause (\href{https://constitutioncenter.org/the-constitution/supreme-court-case-library/dobbs-v-jackson-womens-health-organization}{National Constitution Center}). Invoking \textit{Washington v. Glucksberg}, he insisted that courts must guard against confusing constitutional ``liberty'' with their own policy preferences by anchoring unenumerated rights in history and tradition. Surveying English common law from Bracton, Coke, Hale, and Blackstone through nineteenth-century American statutes, Alito concluded that abortion had been criminalized at common law at least after quickening and that by 1868, when the Fourteenth Amendment was ratified, three-quarters of the States criminalized abortion at all stages of pregnancy (\href{https://constitutioncenter.org/the-constitution/supreme-court-case-library/dobbs-v-jackson-womens-health-organization}{National Constitution Center}). Until shortly before \textit{Roe}, no federal or state court, constitutional provision, or treatise had recognized a constitutional right to abortion.

Second, Alito sought to wall off other substantive due process precedents---\textit{Griswold} (contraception), \textit{Loving} (interracial marriage), \textit{Lawrence} (same-sex intimacy), and \textit{Obergefell} (same-sex marriage)---by arguing that abortion is ``critically different'' because it destroys ``fetal life'' or an ``unborn human being'' (\href{https://constitutioncenter.org/the-constitution/supreme-court-case-library/dobbs-v-jackson-womens-health-organization}{National Constitution Center}).

Third, Alito applied five \textit{stare decisis} factors to justify overruling \textit{Roe} and \textit{Casey}: (1) the nature of the Court's error (usurping democratic authority on a profound moral question); (2) the quality of reasoning (\textit{Roe}'s ``exceptionally weak'' history and legislative-style trimester framework, and \textit{Casey}'s unexplained viability line); (3) unworkability (persistent confusion over \textit{Casey}'s ``undue burden'' test); (4) distortion of other legal doctrines (such as third-party standing, facial challenges, and severability); and (5) absence of concrete reliance interests of the tangible sort found in property or contract law, dismissing \textit{Casey}'s broader societal reliance rationale as an unquantifiable judgment best left to legislatures (\href{https://www.oyez.org/cases/2021/19-1392}{Oyez}; \href{https://www.law.cornell.edu/supremecourt/text/19-1392}{Cornell LII}). Applying rational-basis review, the Court upheld Mississippi's Gestational Age Act as rationally related to legitimate interests in protecting prenatal life and maternal health.

\subsection{Concurring Opinion (Justice Thomas)}

Justice Thomas joined the majority opinion in full but wrote separately to attack the entire doctrine of substantive due process as an oxymoron (\href{https://constitutioncenter.org/the-constitution/supreme-court-case-library/dobbs-v-jackson-womens-health-organization}{National Constitution Center}). Relying on historical evidence that ``due process of law'' guarantees only procedural regularity, Thomas argued that the Due Process Clause never forbids government from infringing substantive rights so long as proper process is provided. Consequently, he urged that ``in future cases, we should reconsider all of this Court's substantive due process precedents, including \textit{Griswold}, \textit{Lawrence}, and \textit{Obergefell},'' calling them ``demonstrably erroneous'' and suggesting that any valid unenumerated rights must instead be evaluated under the Fourteenth Amendment's Privileges or Immunities Clause (\href{https://constitutioncenter.org/the-constitution/supreme-court-case-library/dobbs-v-jackson-womens-health-organization}{National Constitution Center}).

\subsection{Concurring Opinion (Justice Kavanaugh)}

Justice Kavanaugh also joined the majority in full, framing the decision around constitutional neutrality and democratic federalism (\href{https://ballotpedia.org/Dobbs_v._Jackson_Women\%E2\%80\%99s_Health_Organization}{Ballotpedia}). Emphasizing that the interests on both sides of the abortion debate are ``extraordinarily weighty,'' he contended that because the Constitution is silent on abortion, it is ``neither pro-life nor pro-choice'' and leaves the issue to voters and legislators in the States or Congress (\href{https://constitutioncenter.org/the-constitution/supreme-court-case-library/dobbs-v-jackson-womens-health-organization}{National Constitution Center}). Kavanaugh explicitly sought to cabin the opinion's reach, reiterating that overruling \textit{Roe} ``does not threaten or cast doubt on'' \textit{Griswold}, \textit{Eisenstadt}, \textit{Loving}, or \textit{Obergefell}, and volunteering that a State may neither bar a resident from traveling to another State to obtain an abortion (protected by the constitutional right to interstate travel) nor retroactively punish pre-\textit{Dobbs} abortions (barred by the Due Process and Ex Post Facto Clauses) (\href{https://constitutioncenter.org/the-constitution/supreme-court-case-library/dobbs-v-jackson-womens-health-organization}{National Constitution Center}).

\subsection{Concurrence in the Judgment (Chief Justice Roberts)}

Chief Justice Roberts agreed that \textit{Roe} and \textit{Casey}'s viability line should be discarded and Mississippi's 15-week ban upheld, but refused to join the majority in overruling \textit{Roe} and \textit{Casey} ``all the way down to the studs'' (\href{https://constitutioncenter.org/the-constitution/supreme-court-case-library/dobbs-v-jackson-womens-health-organization}{National Constitution Center}). Invoking judicial minimalism---``If it is not necessary to decide more to dispose of a case, then it is necessary not to decide more''---Roberts argued that the Court should have answered only the narrow question on which it granted certiorari (\href{https://constitutioncenter.org/the-constitution/supreme-court-case-library/dobbs-v-jackson-womens-health-organization}{National Constitution Center}). Because the vast majority of abortions occur in the first trimester and a 15-week window provides a ``reasonable opportunity to choose,'' Roberts would have upheld Mississippi's statute while leaving for a future case whether total bans from conception are constitutional.

\subsection{Joint Dissent (Justices Breyer, Sotomayor, and Kagan)}

In a rare jointly authored dissent, Justices Breyer, Sotomayor, and Kagan accused the majority of discarding a 50-year constitutional balance that protected women's bodily autonomy and equal citizenship (\href{https://constitutioncenter.org/the-constitution/supreme-court-case-library/dobbs-v-jackson-womens-health-organization}{National Constitution Center}). They attacked the majority's ``1868 test'' on both historical and methodological grounds: first, early common law did not criminalize pre-quickening abortion; second, freezing the Fourteenth Amendment's meaning to the specific practices of 1868---when women could not vote and were excluded from ``We the People''---consigns women to ``second-class citizenship'' and defies the Framers' choice of open-ended terms like ``liberty'' and ``equality'' (\href{https://constitutioncenter.org/the-constitution/supreme-court-case-library/dobbs-v-jackson-womens-health-organization}{National Constitution Center}).

The dissenters rejected the majority's ``Scout's honor'' assurance that other substantive due process rights are safe, warning that \textit{Griswold}, \textit{Loving}, \textit{Lawrence}, and \textit{Obergefell} all fail the majority's 1868 historical test and stand as part of the same ``constitutional fabric'' (likening the majority's effort to pull out \textit{Roe} to claiming a ``Jenga tower simply will not collapse'') (\href{https://constitutioncenter.org/the-constitution/supreme-court-case-library/dobbs-v-jackson-womens-health-organization}{National Constitution Center}). Finally, they argued that none of the traditional \textit{stare decisis} justifications for overruling precedent---changed facts, doctrinal erosion, or lack of reliance---existed here, concluding that the Court reversed course solely ``because the composition of this Court has changed,'' thereby threatening the Court's legitimacy and the rule of law (\href{https://constitutioncenter.org/the-constitution/supreme-court-case-library/dobbs-v-jackson-womens-health-organization}{National Constitution Center}).

\section{Examples: Future Applications}

\begin{enumerate}
  \item \textbf{Same-Side Application (State Six-Week or Total Abortion Ban):} A State enacts a statute prohibiting all abortions from fertilization onward, with an exception only to save the pregnant woman's life. Under \textit{Dobbs}, a federal Fourteenth Amendment substantive due process challenge will fail because there is no fundamental right to abortion, and the State's asserted interest in protecting unborn life from conception easily satisfies rational-basis review.
  \item \textbf{Same-Side Application (State Restrictions on Medication Abortion Telehealth):} A State prohibits physicians from prescribing abortion-inducing medication via telemedicine and mandates in-person dispensing. Applying \textit{Dobbs}'s rational-basis standard, a federal court would uphold the regulation against a substantive due process challenge because the legislature could rationally conclude that in-person supervision protects maternal health and safety, regardless of empirical disputes over safety data.
  \item \textbf{Opposite-Side Application (State Law Criminalizing Out-of-State Abortion Travel):} A State enacts a statute making it a felony for a resident to travel to another State where abortion is legal to obtain the procedure. Although \textit{Dobbs} eliminates a federal right to abortion within the banning State, a challenge under the constitutional right to interstate travel (and Dormant Commerce Clause / extraterritoriality principles) would likely succeed, as Justice Kavanaugh's concurrence explicitly signaled that a State may not bar residents from traveling across state lines for lawful conduct elsewhere.
  \item \textbf{Opposite-Side Application (State Ban on Married Couples Using Contraception):} A State enacts a law prohibiting the sale or use of oral contraceptives by married couples, arguing that contraception was restricted in the late nineteenth century. Despite Justice Thomas's concurrence, a federal court bound by the \textit{Dobbs} majority opinion would invalidate the law under \textit{Griswold v. Connecticut}, because the \textit{Dobbs} majority expressly distinguished abortion as uniquely involving the destruction of ``potential life'' and declined to disturb non-abortion privacy precedents.
  \item \textbf{Fence-Sitter Application (State Law Restricting In Vitro Fertilization Embryo Destruction):} A State enacts a statute classifying cryopreserved embryos created through in vitro fertilization (``IVF'') as ``unborn human beings'' and prohibiting their disposal or destruction, effectively halting many IVF clinics. This is a genuine fence-sitter: on one hand, \textit{Dobbs}'s history-and-tradition test and deference to state definitions of ``potential life'' strongly favor rational-basis review and upholding the law (as IVF did not exist in 1868); on the other hand, challengers could invoke procreative liberty under \textit{Skinner v. Oklahoma} and \textit{Griswold} (which \textit{Dobbs} left intact) and argue that IVF seeks to \textit{create} rather than terminate life, creating intense doctrinal tension between \textit{Dobbs}'s historical methodology and its stated limitation to abortion.
\end{enumerate}

\section{Critique}

\textbf{Progressive and Doctrinal Criticism:} Critics from progressive, feminist, and living-constitutionalist perspectives contend that \textit{Dobbs} suffers from severe methodological and moral flaws (\href{https://reproductiverights.org/wp-content/uploads/2023/03/Legal-Analysis-What-Dobbs-Got-Wrong-3.15.23.pdf}{Center for Reproductive Rights}; \href{https://www.americanbar.org/groups/communications_law/resources/magazine/2023-summer/how-supreme-courts-decision-dobbs-threatens-privacy-free-speech-rights/}{American Bar Association}). First, they argue that freezing Fourteenth Amendment ``liberty'' to 1868 legal norms perpetuates the historical disenfranchisement of women, who were barred from voting and lawmaking when those nineteenth-century abortion bans were enacted. Second, historians have challenged Alito's reading of the common law, emphasizing that pre-quickening abortion was generally not an indictable offense at common law and that nineteenth-century statutory bans were driven in significant part by physicians' professionalizing campaigns and nativist anxieties rather than universal moral consensus. Third, critics highlight an internal logical contradiction: the majority's \textit{Glucksberg} methodology cannot coherently spare \textit{Griswold}, \textit{Lawrence}, or \textit{Obergefell}, rendering the majority's reassurance that abortion is uniquely different an ad hoc carve-out rather than a principled doctrinal boundary. Fourth, scholars criticize \textit{Dobbs}'s dismissal of \textit{Casey}'s reliance analysis as blind to the lived economic, educational, and bodily reliance of generations of women, and warn that abandoning 50-year-old precedent immediately after personnel changes erodes institutional legitimacy.

\textbf{Originalist and Textualist Perspectives:} Conservative originalist and textualist scholars largely praise \textit{Dobbs} for correcting what even many liberal jurists once conceded was \textit{Roe}'s untethered, legislative-style reasoning, and for restoring democratic self-governance on a deeply contested moral question (\href{https://journals.law.harvard.edu/jlpp/wp-content/uploads/sites/90/2023/10/24_46_S_Walsh.pdf}{Harvard Journal of Law \& Public Policy}). Yet \textit{Dobbs} has also drawn sharp criticism from within originalism itself (\href{https://ir.lawnet.fordham.edu/cgi/viewcontent.cgi?article=2923&context=ulj}{Fordham Urban Law Journal}). Strict textualists echoing Justice Thomas fault Alito for retaining the doctrine of substantive due process and the \textit{Glucksberg} test at all, rather than grounding unenumerated rights analysis in the original public meaning of the Privileges or Immunities Clause. Conversely, some originalist scholars argue that Alito's ``history and tradition'' inquiry conflated original public meaning with original expected applications, relying on state statutory counts in 1868 without rigorously analyzing the semantic scope of ``liberty'' or whether unborn fetuses were understood as ``persons'' under the Fourteenth Amendment (a question the majority expressly declined to decide). Meanwhile, judicial minimalists aligned with Chief Justice Roberts argue that overruling \textit{Roe} and \textit{Casey} in full when a narrower 15-week ruling would have resolved the case breached Burkean judicial restraint.

\section{Key Quotations}

\begin{quote}
``We hold that \textit{Roe} and \textit{Casey} must be overruled. The Constitution makes no reference to abortion, and no such right is implicitly protected by any constitutional provision, including the one on which the defenders of \textit{Roe} and \textit{Casey} now chiefly rely---the Due Process Clause of the Fourteenth Amendment.''
\end{quote}
\noindent --- Justice Samuel A. Alito, Jr., Opinion of the Court (\href{https://constitutioncenter.org/the-constitution/supreme-court-case-library/dobbs-v-jackson-womens-health-organization}{National Constitution Center}).

\begin{quote}
``\textit{Roe} was egregiously wrong from the start. Its reasoning was exceptionally weak, and the decision has had damaging consequences. And far from bringing about a national settlement of the abortion issue, \textit{Roe} and \textit{Casey} have enflamed debate and deepened division. It is time to heed the Constitution and return the issue of abortion to the people's elected representatives.''
\end{quote}
\noindent --- Justice Samuel A. Alito, Jr., Opinion of the Court (\href{https://www.law.cornell.edu/supremecourt/text/19-1392}{Cornell LII}).

\begin{quote}
``On the question of abortion, the Constitution is therefore neither pro-life nor pro-choice. The Constitution is neutral and leaves the issue for the people and their elected representatives to resolve through the democratic process in the States or Congress---like the numerous other difficult questions of American social and economic policy that the Constitution does not address.''
\end{quote}
\noindent --- Justice Brett M. Kavanaugh, Concurring Opinion (\href{https://ballotpedia.org/Dobbs_v._Jackson_Women\%E2\%80\%99s_Health_Organization}{Ballotpedia}).

\begin{quote}
``If it is not necessary to decide more to dispose of a case, then it is necessary not to decide more.''
\end{quote}
\noindent --- Chief Justice John G. Roberts, Jr., Concurring in the Judgment (\href{https://constitutioncenter.org/the-constitution/supreme-court-case-library/dobbs-v-jackson-womens-health-organization}{National Constitution Center}).

\begin{quote}
``With sorrow---for this Court, but more, for the many millions of American women who have today lost a fundamental constitutional protection---we dissent.''
\end{quote}
\noindent --- Justices Stephen G. Breyer, Sonia Sotomayor, and Elena Kagan, Joint Dissent (\href{https://constitutioncenter.org/the-constitution/supreme-court-case-library/dobbs-v-jackson-womens-health-organization}{National Constitution Center}).


\end{document}
